\markright{\mititulo \hspace{5cm} \thepage}
\pagestyle{fancy}

\section*{Resumen}
\addcontentsline{toc}{section}{Resumen}

\noindent En esta práctica se analizaron dos evidencias clásicas del heliocentrismo
mediante efemérides de \textit{NASA/JPL Horizons} consultadas con
\texttt{astroquery}: (i) la relación entre fase y tamaño angular de Venus y
(ii) la retrogradación aparente de Marte.
Para Venus se estudió un intervalo de 608 días (2022-06-01 a 2024-03-01), que
cubre un período sinódico completo ($P_{\text{syn}}\approx 583.93$ días), y se
observó una anti-correlación marcada entre fracción iluminada y diámetro
aparente: cuando Venus está cerca de la Tierra es grande y poco iluminado, y
cuando está lejos es pequeño y casi lleno.
Para Marte se examinó el periodo 2026-09-01 a 2027-09-01 con paso de 10 días,
identificando una oposición el 2027-02-18 (elongación
$174.96^\circ$, diámetro $13.81''$, distancia $0.6784$ UA) y un intervalo
retrógrado del 2027-01-09 al 2027-03-30 (80 días).
Los resultados reproducen cuantitativamente los fenómenos históricos usados para
rechazar el modelo geocéntrico y confirman que las variaciones observadas se
explican de forma natural por la geometría orbital heliocéntrica.

\vspace{1cm}\textit{Palabras clave:} heliocentrismo, Venus, Marte, retrogradación,
tamaño angular, fases planetarias, JPL Horizons, astroquery

\newpage

\section*{Abstract}
\addcontentsline{toc}{section}{Abstract}

\noindent This report analyzes two classical heliocentric evidences using
\textit{NASA/JPL Horizons} ephemerides queried through \texttt{astroquery}:
(i) the relation between Venus phase and apparent angular diameter and
(ii) the apparent retrograde motion of Mars.
For Venus, a 608-day interval (2022-06-01 to 2024-03-01) was studied,
covering one synodic cycle ($P_{\text{syn}}\approx 583.93$ days), and showing
a strong anti-correlation between illuminated fraction and angular diameter:
when Venus is near Earth it appears large and weakly illuminated, while when
it is far it appears small and almost full.
For Mars, the interval 2026-09-01 to 2027-09-01 (10-day step) was analyzed,
identifying an opposition on 2027-02-18 (solar elongation
$174.96^\circ$, diameter $13.81''$, distance $0.6784$ AU) and a retrograde
interval from 2027-01-09 to 2027-03-30 (80 days).
The results quantitatively reproduce the historical observations that falsified
the geocentric model and confirm that the observed variations are naturally
explained by heliocentric orbital geometry.

\vspace{1cm}\textit{Keywords:} heliocentrism, Venus, Mars, retrograde motion,
angular diameter, planetary phases, JPL Horizons, astroquery

\newpage
